% ===== Document Preamble =====
\documentclass[11pt]{article}
\usepackage[margin=0.5in]{geometry}
\usepackage{amsmath,amssymb}
\usepackage{xcolor}
\usepackage{enumitem}
\usepackage{needspace}
\usepackage{comment}

% ===== Colors =====
\definecolor{AnsRed}{RGB}{170,0,0}
\definecolor{WrongBlue}{RGB}{0,0,160}

% ===== Choice list styling (tight, consistent) =====
\setlist[enumerate]{%
  leftmargin=1.5em,
  itemsep=0.25em,
  topsep=0.25em,
  label=\textbf{\Alph*.}
}

% ===== Style Wrappers =====
\newcommand{\QuestionStyle}[1]{{\large\bfseries\boldmath #1}}
\newcommand{\AnswerStyle}[1]{{\large\bfseries\boldmath #1}}
\newcommand{\ReasonStyle}[1]{{\normalsize #1}}

% ===== Show / Hide Reasoning Toggle =====
\newif\ifshowreason
\showreasontrue      % Instructor version
% \showreasonfalse   % Student version

% ===== Choice Item Macros (Answer + Reasoning) =====
% Usage:
%   \Correct{Answer}{Reasoning}
%   \Wrong{Answer}{Reasoning}

\newcommand{\Correct}[2]{%
  \item \AnswerStyle{\textcolor{AnsRed}{#1}}%
  \ifshowreason\\[-0.2em]
    \ReasonStyle{\textcolor{AnsRed}{#2}}%
  \fi
}

\newcommand{\Wrong}[2]{%
  \item \AnswerStyle{\textcolor{WrongBlue}{#1}}%
  \ifshowreason\\[-0.2em]
    \ReasonStyle{\textcolor{WrongBlue}{#2}}%
  \fi
}

% ===== Question Block Macro =====
% Usage: \QBlock{Q1}{stem...}{<items>}

\newcommand{\QBlock}[3]{%
  \Needspace{8\baselineskip}%
  \noindent
  \begin{minipage}{\linewidth}
    \QuestionStyle{#1.~#2}

    \vspace{0.25em}
    \begin{enumerate}
      #3
    \end{enumerate}

    \vspace{0.25em}
    \hrule height 0.4pt
  \end{minipage}
  \vspace{0.75em}
}

% ===== Document Start =====

\begin{document}

\begin{center}
{\huge \textbf{Isomorphic Vector and Kinematics Questions}}\\[0.25em]
{\Large Position vectors, displacement, velocity, and acceleration}\\[0.25em]
{\large Multiple choice questions generated below. Answers in RED and incorrect answers in BLUE.}
\end{center}

\vspace{0.75em}

\QBlock{Q1}{A search-and-rescue drone reports its position relative to the command post located at the origin. At a certain moment its position vector is $\vec{r}= -7\hat{i}+9\hat{j}$, where distances are measured in kilometers. Which statement is correct?}{
  \Wrong{The drone is $16\ \text{km}$ from the command post.}{This corresponds to adding component magnitudes ($7+9$) instead of using the Euclidean distance formula.}
  \Correct{The drone is approximately $11.4\ \text{km}$ from the command post.}{The distance from the origin is the magnitude of the position vector: $|\vec{r}|=\sqrt{(-7)^2+9^2}=\sqrt{130}\approx 11.4\ \text{km}$. This uses the Pythagorean theorem for perpendicular components in the $x$- and $y$-directions.}
  \Wrong{The drone is $9\ \text{km}$ from the command post.}{This incorrectly treats the $y$-component ($9$ km) as the full distance from the origin.}
  \Wrong{The drone is $7\ \text{km}$ from the command post.}{This incorrectly treats the $x$-component magnitude ($7$ km) as the full distance from the origin.}
}

\QBlock{Q2}{A cart moves along the $x$-axis with displacement given by $\vec{s}(t)=\Big(\left(0.80\,\mathrm{m/s^3}\right)t^3+\left(4.5\,\mathrm{m/s^2}\right)t^2\Big)\hat{i}$. Which statement is correct?}{
  \Correct{The magnitude of the acceleration at $t=0\ \text{s}$ is $9.0\ \text{m/s}^2$.}{Differentiate twice: $\vec{v}(t)=\dfrac{d\vec{s}}{dt}=(2.40t^2+9.0t)\hat{i}$ and $\vec{a}(t)=\dfrac{d\vec{v}}{dt}=(4.8t+9.0)\hat{i}$. Evaluating at $t=0$ gives $|\vec{a}(0)|=9.0\ \mathrm{m/s^2}$.}
  \Wrong{This motion has constant acceleration.}{Taking the second derivative gives $a(t)=4.8t+9.0$, which depends on $t$, so the acceleration is not constant.}
  \Wrong{The magnitude of the acceleration is $4.5\ \text{m/s}^2$.}{The coefficient $4.5\,\mathrm{m/s^2}$ appears in the displacement function, but acceleration comes from the second derivative, not from reading off this coefficient.}
  \Wrong{The magnitude of the acceleration at $t=0\ \text{s}$ is $0\ \text{m/s}^2$.}{A common error is assuming an object starts with zero acceleration; here $a(0)=2\times 4.5=9.0\ \mathrm{m/s^2}$ from $\vec{a}(t)=\dfrac{d^2\vec{s}}{dt^2}$.}
}

\QBlock{Q3}{A particle moves in the $xy$-plane from point A at $\vec{r}_A=-3\hat{i}+4\hat{j}$ to point B at $\vec{r}_B=5\hat{i}-8\hat{j}$ in $4.0\ \text{s}$ (coordinates in meters). Which statement is correct?}{
  \Wrong{The straight-line distance between A and B is $20\ \text{m}$.}{This comes from adding component changes ($|\Delta x|+|\Delta y|=8+12$) instead of using $\sqrt{(\Delta x)^2+(\Delta y)^2}$.}
  \Wrong{The speed at $t=0$ must be $3.6\ \text{m/s}$.}{Computing $|\Delta \vec{r}|/\Delta t$ gives the magnitude of average velocity, not the instantaneous speed at $t=0$; many motions can share the same endpoints and time.}
  \Correct{The straight-line distance between A and B is approximately $14.4\ \text{m}$.}{Compute displacement: $\Delta\vec{r}=\vec{r}_B-\vec{r}_A=(5-(-3))\hat{i}+(-8-4)\hat{j}=8\hat{i}-12\hat{j}$. The straight-line separation is $|\Delta\vec{r}|=\sqrt{8^2+(-12)^2}=\sqrt{208}\approx 14.4\ \text{m}$.}
  \Wrong{The acceleration at $t=0$ must be $0.90\ \text{m/s}^2$.}{Acceleration at a particular time cannot be uniquely determined from only initial/final positions and elapsed time without extra information (e.g., full trajectory or a constant-acceleration constraint plus more data).}
}

\end{document}
